\documentclass[10pt,a4paper]{amsart}
\usepackage[margin=19mm]{geometry}
\usepackage{amsmath,amssymb,mathtools}
\usepackage[scr=rsfs,cal=euler]{mathalfa}
\usepackage{xfrac}
\usepackage[unicode,hidelinks,bookmarks=false]{hyperref}
\newcommand{\ie}{i.e.\ }
\newcommand{\cf}{cf.\ }
\newcommand{\resp}{respectively\ }
\newcommand{\Subsection}{\subsection}
\providecommand{\nobreakdash}{\nobreak\hbox{-}}
\setlength{\emergencystretch}{2em}
\multlinegap0pt
\usepackage{bm}
\usepackage{bbm}
\usepackage{dsfont}
\usepackage{xspace}
\usepackage{cancel}
\usepackage{mathtools}
\usepackage{amsthm}
\makeatletter
\@ifundefined{theorem}{\newtheorem{theorem}{Theorem}}{}
\makeatother
\title[On Targeted Complexity of Discrete Motion]{On Targeted Complexity of Discrete Motion}
\author{Ameneh Babaee, Hanieh Mirebrahimi, Hamid Torabi, Soheila Fahimi}
\date{Source: arXiv:2508.07052v2 (CC BY 4.0)}
\begin{document}
\maketitle

\noindent\emph{Source and licence.} On Targeted Complexity of Discrete Motion. Version arXiv:2508.07052v2 (CC BY 4.0), distributed under the Creative Commons Attribution 4.0 International licence, as recorded in the arXiv licence metadata for this exact version and shown on the abstract page https://arxiv.org/abs/2508.07052v2. Full rights notice: licenses/arxiv-2508.07052-CC-BY-4.0.txt.\par\smallskip

\noindent\textit{Editorial scope.} Counted unit: the Theorem whose source label is \texttt{eqnw1} in the article (source0.tex lines 1360--1395, source label \texttt{eqnw1}), delivered together with the article's own complete original proof of it (source0.tex lines 1397--1435, one \texttt{proof} environment of 39 lines). No mathematics is written, repaired or re-derived: statement and proof are the article's own text line for line. Required context retained uncounted: none. Every definition, display and label of those lines is retained unchanged. Omitted from this package and not redistributed: 1--1359, 1396--1396, 1436--1599. Nothing inside a retained excerpt is omitted or altered: every retained body is delivered exactly as the article's lines are written, so no omission receipt of any kind accompanies this package. Citations: the retained excerpts carry no citation command at all, so no bibliography entry of the article is transcribed and no reference is redistributed. Reference adaptation: none was needed and none was made; every reference inside the delivered unit resolves inside it, so no unresolved reference or citation is produced. Printed slips kept literal: no printed word, symbol, hypothesis or number of the counted statement or of its proof was altered. Layout and class: the article's own class is \texttt{hjm1} with packages including xy, subfig, amsthm, amsfonts, amssymb, amscd, amsmath, enumerate, verbatim, calc, graphicx, geometry, tikz; the wrapper is the lane's pinned \texttt{amsart} editorial wrapper at 10pt on A4, which restates the theorem-like environments the retained text needs and adds the article's own macro definitions that the retained text needs (none), with extra wrapper packages (bm, bbm, dsfont, xspace, cancel, mathtools). The delivered numbering is the unit's own. Assets: no retained span contains a figure environment, a graphics-inclusion command, a PDF-page inclusion, a table or a TikZ picture; no image file is copied into this package, the package declares no asset dependency and no image file is redistributed with it. Licence: version arXiv:2508.07052v2 of this article is distributed under the Creative Commons Attribution 4.0 International licence, as recorded in the arXiv licence metadata for this exact version and shown on the abstract page https://arxiv.org/abs/2508.07052v2. A full rights notice is delivered at licenses/arxiv-2508.07052-CC-BY-4.0.txt. Characters: every retained excerpt is pure ASCII, so the delivered engine, XeLaTeX with its default Latin Modern text font, reports no missing character.
\begin{theorem}
Let $n$ be a natural number, and $K$ be a finite $[\frac{n}{2}]$-connected simplicial complex.
\begin{enumerate}
\item
If  $v_* \in K$, then
\begin{align}
& scat_{n } (K) \le TC_{n} (K,v_*) \le scat_{[\frac{n}{2}]} (K), \label{eqnw1}
%\\
%& scat_{2m +1} (K) \le TC_{2m +1} (K,v_*) \le scat_{m} (K).\label{eqnw2}
\end{align}
% $scat_m(K)=TC_m(K, v_*)$.
%\end{theorem}


%\begin{theorem}
\item
If $L\subseteq K$ be a non-empty subcomplex of $K$, then 
\begin{align}
& scat_{n}(K)\leq TC_{n}(K,L) \le TC_{n}(K) \leq scat_{[\frac{n}{2}]}(K\prod K), \label{eqnw3}
%\\
%& scat_{2m+1}(K)\leq TC_{2m +1}(K,L) \le TC_{2m +1}(K) \leq scat_m(K\prod K). \label{eqnw4}
\end{align}
%\end{theorem}

%\item
%Let $K$ be a finite connected complex and $L$ be a $m$-categorical subcomplex of $K$. Then 
%\begin{align}
% & scat_{2m} (K) \le TC_{2m}(K,L) \le scat_m(K), \label{eqnw5}
%\\
%& scat_{2m+1} (K) \le TC_{2m+1}(K,L) \le scat_m(K). \label{eqnw6}
%\end{align}

%\item
%$TC_m(K,L)=0$ if and only if $K$ is $m$-strongly collapsible.
\end{enumerate}
\end{theorem}

\begin{proof}
\begin{enumerate}
\item
Let $n$ be an even number, $m = \frac{n}{2}$ and  $scat_{m} (K) = r$. Then there exist $m$-categorical subcomplexes $\Lambda _0,...,\Lambda _r$ of $K$ and $m$-homotopies $F_i :\iota_{\Lambda_i} \simeq c_{v}$ such that $\omega_1 \circ \lambda _i=\iota_{\Lambda _i}$ for some $v \in K$ and for all $0\leq i\leq r$. 
Define $\Omega _i =\Lambda _i \prod \lbrace v_* \rbrace $. It is clear that $\lbrace \Omega _i \rbrace _0 ^r$ covers $K\prod \lbrace  v_*\rbrace$. Put $\sigma _i:\Omega _i\to P_{n}(K,v_*)$ by $\sigma _i (v,v_*)=F_i(v, -) * \gamma_i$ where $\gamma_i$ is the $m$-path from $v$ to $v_*$. Therefore $TC_{n}(K, v_*) \le scat_m(K) = scat_{[\frac{n}{2}]} (K)$.


Now let $TC_{n} (K, v_*) = r$. Then there exist subcomplexes $\Omega_0, \ldots, \Omega_r$ of $K \prod \{v_*\}$ and simplicial maps $\sigma_i: \Omega_i \to P_{n} (K,v_*)$ such that $\pi_{n} \circ \sigma_i = \iota_{\Omega_i}$. Define $\Lambda_i$ as the projection of $\Omega_i$ on the first component and  also define $F_i: \lambda_i \to P_{n} K$ as composition $F_i: \Lambda_i \hookrightarrow \Omega_i \to^{\sigma_i} P_{n}(K, v_*) \hookrightarrow P_{n} K$.  Since $\alpha \circ \sigma_i = \iota_{\Omega_i}$, we have $\alpha \circ F_i = \iota_{\Lambda_i}$ and also $\omega \circ F_i = c_{v_*}$. Therefore $\Lambda_i$ is $n$-categorical and then $scat_{n} (K) \le TC_{n} (K, v_*)$. 

For the odd index, let $m = [\frac{n}{2}]$. Then $n = 2m +1$ and we prove the sequence of inequalities \eqref{eqnw1} in odd cases, 
by using this fact that $TC_{2m +1} \le TC_{2m}$. This yields to the right inequality and the left inequality is obtained similar to the left inequality of 
\eqref{eqnw1}.
%$F_i(v, m) \neq v_0$???
%
%%
%% We have $p_0 \circ \sigma _i(v,v_0)=(\alpha (\sigma _i(v,v_0)),\omega(\sigma _i(v,v_0)))=(\alpha \overline{(\lambda _i(v))},\omega (\overline{\lambda _i(v))})=(v,v_0)$.
%%Then $secat(p_0)\leq secat(\omega)$. 
%Now suppose that $TC_m (K, v_0)=n$. Thus there exist subcomplexes $\Omega _0,...,\Omega _n$ of $K\prod \lbrace v_0 \rbrace$ and simplicial maps $\sigma _i:\Omega _i\to P_m(K,v_0)$ such that $\pi_m \circ \sigma _i=i_{\Omega _i}$ for all $0\leq i\leq m$. We define $\Lambda _i$ as the projection of $\Omega _i$ on the first component and $\lambda _i:\Lambda _i\to P_m(K, v_0)$ by $\lambda _i(v)=\overline{\sigma _i(v,v_0)}$. Hence we have $\omega \circ \lambda _i(v)=\omega \overline{(\sigma _i (v,v_0))}=v$ and  then $secat(\omega)\leq secat(p_0)$. Therefore $secat(\omega)= secat(p_o)$.
%%\end{proof}
%
%
%Now Let $K$ be a connected finite complex. Since $K$ is finite, by Theorem \ref{i}, $TC(K,v_0)=secat (p_0)$.
%Also since $K$ is connected and finite, by Theorem \ref{th2.13n}, $scat (K) = secat (\omega_1)$. Thus $TC(K,v _0)=scat(K)$.

\item
Let $n$ be an even number, $m = \frac{n}{2}$ and   $scat_m(K \prod K) = r$. Then there exist $m$-categorical subcomplexes $\Lambda_0, \ldots, \Lambda_r$ of $K \prod K$. Define $\sigma_i (v,v')= p_1 \circ F_i (v, v') * p_2 \circ \overline{F_i (v,v')}$ where $F_i: \Lambda_i \to P_m( K \prod K)$ is the $m$-homotopy between $\iota_{\Lambda_i}$ and the constant map         and $p_1$ and $p_2$ are the projection maps on the first and second components.  
Hence $TC_{n} (K) \le scat_m(K) = scat_{[\frac{n}{2}]} (K)$. 
Now let $TC_{n} (K) =r$. Then similar to the item 1, one can prove that $scat_{n} (K) \le TC_{n} (K,L)$.

Similar to the item (1), for the odd index, again let $m = [\frac{n}{2}]$. Then $n = 2m +1$ and  the sequence of inequalities \eqref{eqnw3} in odd cases, is proved by the inequality  $TC_{2m +1} \le TC_{2m}$.


%\item
%We know that $scat_m(K) \le TC_m(K,L)$. Now let $scat_m(K) =n$. Then there exist $m$-categorical subcomplexes $\Lambda_0, \ldots, \Lambda_n$ of $K$. Put $\Omega_i = \Lambda_i \prod L$ and $\sigma_i = F_i \prod G$ where $F_i: \iota_{\Lambda_i} \simeq c_v$ and $G:\iota_{L} \simeq c_{v'}$. 


 
\end{enumerate}
\end{proof}

\end{document}
